
In last milestone we assumed a common interaction strength across all edges $\beta_{ij}=\beta.$ and we also assumed a linear model $\sum_{j:j\sim i}\alpha_{ij}=2X_i^\top\theta$.
In this setting, however, the graph has three different edge types. It is therefore too restrictive to use a single interaction parameter. The preceding model treats all edges of the same relation as equally
informative. This is a strong assumption in fraud-detection graphs. An edge only indicates that two nodes share a certain relation; it does not necessarily mean that the two nodes should influence each other equally. Two connected nodes that are very similar in their features should plausibly exert more influence on each other than two connected nodes whose features are very different. Therefore, we replace the unweighted interaction by a similarity-weighted interaction.

For an edge $(i,j)\in E_r$, define a feature-similarity weight $w_{ij}^{(r)}$
In our implementation, we use the RBF kernel
$$
    w_{ij}^{(r)}
    =
    \exp\left(
    -\frac{\|X_i-X_j\|_2^2}{2\tau_r^2}
    \right),
$$
where $\tau_r>0$ is a bandwidth hyper-parameter for relation $r$. This gives high weight to feature-similar neighbors and low weight to feature-dissimilar neighbors.

This leads to the weighted interaction model
$$
    \mathbb P_{\beta,\theta}(\sigma\mid X,G_1,G_2,G_3)
    \propto
    \exp\left\{
    \sum_{r=1}^3
    \beta_r
    \sum_{(i,j)\in E_r}
    w_{ij}^{(r)}\sigma_i\sigma_j
    +
    2\sum_{i=1}^n\sigma_i X_i^\top\theta
    \right\}.
$$

There is one further issue: node degrees can vary substantially. If we use the raw weighted sum $\sum_{j:(i,j)\in E_r}w_{ij}^{(r)}\sigma_j,$ then high-degree nodes can have much larger local fields simply because they
have many neighbors. To make the graph statistics comparable across nodes, we normalize using weighted neighbor average,$D_i^{(r)}:=\sum_{k=1}^nA_{ik}^{(r)}w_{ik}^{(r)},$ for $r=1,2,3.$

\[
    \mathbb P_{\beta,\theta}(\sigma\mid X,G_1,G_2,G_3)
    \propto
    \exp\left\{
    \sum_{r=1}^3
    \beta_r
    \sum_{i=1}^n
    \sum_{j=1}^n
    A_{ij}^{(r)}w_{ij}^{(r)}
    \left(
    \frac{1}{D_i^{(r)}}+\frac{1}{D_j^{(r)}}
    \right)
    \sigma_i\sigma_j
    +
    2\sum_{i=1}^n \sigma_i X_i^\top\theta
    \right\}.
\]

Finally, instead of forcing the node-level feature effect to be linear, we also consider a nonlinear feature effect. We replace $X_i^\top\theta$ by a scalar-valued function $b_\theta(X_i).$ Thus the neural weighted multi-relation model is
$$
    \mathbb P_{\beta,\theta}(\sigma\mid X,G_1,G_2,G_3)\propto\exp\left\{\sum_{r=1}^3\beta_r\sum_{(i,j)\in E_r}w_{ij}^{(r)}\sigma_i\sigma_j+2\sum_{i=1}^n\sigma_i b_\theta(X_i)\right\}.
$$

This formulation preserves the original network-regularized logistic idea while adapting it to a multi-relation graph, similarity-weighted edges, degree normalization, and nonlinear covariate effects.